% results_b.tex : 결과 R5-R8(국문판). 영문판 paper/manuscript/en/sections/results_b.tex(2026-10-05 2차 수정판)를
% 문장 단위로 옮겼다. 수치, 인용 키, 그림·보충 표 번호, 자리표시 문자열은 영문판과 같다.

\subsection{라벨이 많은 지역 안의 이득}
\label{sec:results-label-rich}

알래스카 안에서 재보정 Stefan 앵커 위의 잔차 ML은 같은 라벨로 재보정한 Stefan 식보다 라벨 1000개에서 RMSE를 0.54~cm 낮추었다(95\% CI [0.32, 0.69], 블록 등가중 [0.40, 0.80], 14.4~cm에서 13.9~cm, 공통 재표집에서도 같음). 라벨 전량에서 오차 제곱 감소의 약 99\%는 ERA5-Land 기후 격자 셀 사이에 있었다\cite{munozsabater2021_era5_land}(그림~7a,c).
이 시험은 앞선 결과를 본 뒤 설계하고 실행 전에 등록하였다(방법).

라벨 500개와 전량의 감소(0.39~cm와 0.52~cm)는 분할 독립 가정에 의존하며, 원천 계수 Stefan 식과의 차이는 라벨 1000개와 전량에서 확인하지 못했다(보충 표~S7).
레나델타와 캐나다에서는 같은 대비의 차이를 시험한 어느 라벨 수에서도 확인하지 못했다(보충 표~S7).

오차를 기후 격자 셀 사이와 안으로 나누면(대용: 0.5° 블록과 ERA5-Land TDD; 방법; 앞선 결과를 본 뒤 설계, 실행 전 등록), 라벨 전량의 잔차 ML은 재보정 모델보다 격자 사이 RMSE를 1.0~cm 낮추었고(95\% CI [0.6, 1.3]), 격자 안 RMSE는 ±0.5~cm 안에서 동등하였다(−0.01~cm, 95\% CI [−0.03, 0.04]). 토양 온도 격자 정의에서도 같았다(그림~7c).
격자 안 변동은 재보정 Stefan 식 오차 제곱의 72\%를 차지하였고, 잔차 ML은 그 가운데 알래스카에서 1\% 미만, 세 지역의 어느 설정에서도 8\% 미만을 설명하였다.
이득을 격자 규모의 편향 보정으로 보는 것은 사후 서술이다(채점 블록별 오차 변화는 그림~7b).

등록된 격자 안 가설에서 잔차 ML의 격자 안 예측은 라벨 전량에서 층화 지역 평균의 오차를 키웠다(+0.09~cm, 95\% CI [0.02, 0.40]). 토양 온도 정의와 캐나다에서도 같았고, 모두 통계적으로 구별되나 크기는 0.5~cm 미만이었다(보충 표~S10).
전이에서는 격자 안 예측이 관측 변동과 맞는다는 근거가 없었다.
더 따뜻한 블록으로의 외삽에 대한 등록 시험은 판정할 수 없었다(캐나다의 채점 블록 3개; 보충 표~S10).
다음 세 분석은 앞선 결과를 본 뒤 설계하였다(보충 표~S13).
같은 시기의 따뜻한 블록을 쓴 공간 대용(외삽 폭은 $\sqrt{\mathrm{TDD}}$에서 0.97~($^{\circ}$C~d)$^{1/2}$)은 실제 외삽 시험이 아니며, 여기서 잔차 ML은 재보정 Stefan 식보다 라벨 100개와 전량에서 오차가 2.44~cm와 2.41~cm 낮았고(비열등), 직접 ML과 두 모델 사이의 외삽 손실 차이는 확인하지 못했다.
지형, 수체, 수목 피복 입력 10개를 더해도 잔차 ML의 격자 안 오차는 모든 라벨 수에서 ±0.5~cm 안에서 동등하였다(보조 대비, 비맹검 부분 포함). \textcolor{blue}{[PENDING: XE-a, XE-b, satellite inputs (stage 2); registered deadline 8 October 2026, 14:22 KST]}
적층 잔차는 라벨 500개와 1000개에서 재보정 앵커 잔차보다 오차가 0.55~cm와 0.498~cm 낮았고(3지역 가운데 2지역; 뒤의 것은 0.5~cm 미만; CCI v5를 후보로 더하면 유지되지 않음), 라벨 200개와 전량에서는 차이를 확인하지 못했으며, 알래스카에서 적층은 재보정 Stefan 식에 비열등이었다(비맹검 부분 포함).

\subsection{새 라벨의 위치}
\label{sec:results-placement}

캐나다에서 라벨을 0.5° 블록이나 공변량 공간에 흩어 놓으면 무작위 셀 추출보다 RMSE가 전이 라벨 40개에서 2.4~cm(95\% CI [0.5, 3.4], 블록 등가중 [0.8, 2.5], 공변량 최대 최소 거리 추출), 지역 안 라벨 50--100개에서 2.5--3.2~cm 낮아졌다. 레나델타에서는 블록 층화가 라벨 100개에서 RMSE를 1.7~cm 높였다(95\% CI [0.9, 2.4], 블록 등가중 [0.1, 1.6])(그림~5).
모든 대비는 같은 분할에서 잔차 ML($\lambda = 0.25$)을 썼고, 앞선 결과를 본 뒤 설계하고 실행 전에 등록하였다.

전이 라벨 10개에서 어느 전략도 다섯 지역 가운데 오차가 높아진 곳이 없었으나, 이는 비열등 결과가 아니다. 한 가중에서는 레나델타의 블록 층화(+1.9~cm)와 러시아 동부의 공변량 분산(+0.36~cm)이 오차 증가를 보였기 때문이다(그림~5f; 보충 표~S9).

라벨 100개에서 알래스카 최선 전략인 공변량 군집 층화는 레나델타에서 RMSE를 높였다(+0.33~cm, 0.5~cm 미만).
내부 사전 등록한 블록 층화 대 무작위 셀 추출의 지역 평균 대비는 라벨 10개와 40개에서 미결정이었다(2단 재표집; 비맹검 부분 포함; 전이 대비와 같은 라벨 집합; 보충 표~S9).
평가 전에 고정한 배치 절차(방법)도 순차 워크플로 안에서 시험하였다(그림~7e).
같은 지역을 다시 무작위로 나눈 분할에서 이 절차는 셀 무작위 추출로 바뀌었으므로 배치 단계는 기여한 것이 없고, 그 효과는 시험하지 않았다(방법). 이 절차는 알래스카 계열에서 골랐고, 평가 지역의 셀은 배치와 방법 규칙의 설계에 쓰였으며, 평가 분할은 같은 셀을 새로 나눈 것이다.
블록 배분과 공변량 덮임을 합친 배치는 라벨 10개와 40개에서 무작위 배치와의 차이를 확인하지 못했고, 그 네 설정은 블록 또는 공변량 분산 배치와의 비교 16개 가운데 8개에서 비열등(2개는 보정 전에만 유의), 2개에서 오차가 높았으며(0.5~cm 미만), 레나델타를 겨냥한 사후 설계인 설계 가중 계수는 알래스카와 캐나다에서 오차를 키웠다(탐색; 앞선 결과를 본 뒤 설계, 비맹검 부분 포함; 보충 표~S13).
후보는 이미 측정한 셀이므로 현장 후보로 이득이 옮겨진다는 보장이 없다.

\subsection{독립 지역}
\label{sec:results-independent}

얕은 활동층을 가진 약관 확인 새 지역인 러시아 중부(라벨 57개)에서 대상 라벨 없는 직접 ML은 원천 계수 Stefan 식보다 RMSE를 2.7~cm 높였다(95\% CI [2.1, 4.1], 블록 등가중 [3.2, 6.6], 공통 재표집에서도 같음). 이 지역을 포함한 얕은 레짐 5지역 풀에서 라벨 전량의 잔차 ML은 같은 모델보다 RMSE를 3.2~cm 낮추었다(95\% CI [2.2, 4.2], 블록 등가중 [1.5, 3.7])(표~1).
[MISSING: 5지역 풀 대비(라벨 전량 잔차 ML − 원천 계수 Stefan, 라벨 10개 잔차 ML − 재보정 Stefan)의 지역 중앙값, 오차 감소·증가 지역 수, 러시아 서부 제외 값. 새 지역 판정표(약관 확인분 판)의 지역 행에서 수치 생성 스크립트로 만든다]
이 결과에는 교차 환경 표지가 붙는다. 풀이 두 계산 플랫폼에서 맞춘 지역을 합치기 때문이다(방법).
4지역 풀과 5지역 풀의 어느 지역에서도 라벨 0--40개의 직접 ML은 원천 계수 모델보다 오차가 낮지 않았다.
라벨 10개에서 잔차 ML과 재보정 Stefan 식은 5지역 풀에서 ±0.5~cm 안에서 동등하였으므로(−0.13~cm), 이 라벨 수의 4지역 결과는 이 풀에서 유지되지 않았다.
러시아 중부만 보면 라벨 전량의 잔차 ML은 차이를 확인하지 못했다. 두 가중의 판정이 엇갈렸기 때문이다(보충 표~S11).

깊은 활동층을 가진 티베트 고원은 따로 보고하며 독립 지역의 근거로 쓰지 않는다(채점 셀이 원천 고도 범위 밖에 있음; 비맹검 대비).
그곳에서 라벨 10개의 잔차 ML은 수축 재보정보다 RMSE가 20.9~cm 낮았으나, 같은 라벨에 최소제곱으로 맞춘 Stefan 식보다 낫지 않았다(+19.1~cm; 보충 표~S11).
자료 약관이 일부 확인되지 않은 북대서양 지역은 보충 표~S11에만 있다.
\textcolor{blue}{[PENDING: XF, new public regions; registered data deadline 11 October 2026, 14:22 KST; no eligible new macro region found so far]}
\textcolor{blue}{[PENDING: XC-F3, external-family test of the workflow; not tested so far, because placement reduced to random sampling makes its ten-label contrast identical, and no eligible new region; registered sentence after the XF data deadline, 11 October 2026, 14:22 KST]}
티베트 고원에서 지온 유도 보조 행은 라벨 3개와 10개로 재보정한 모델의 전이 오차를 8.14--83.17~cm 줄였다(앞선 결과를 본 뒤 설계, 비맹검 부분 포함). 이 라벨은 직접 측정과 정의가 달라 이 결과를 일반화하지 않으며, 원천 쪽 설계는 시험하지 않았다(약관 근거 미기록; 보충 표~S13).

\subsection{예측 구간}
\label{sec:results-intervals}

대상 라벨이 없을 때 지역을 묶음으로 보정한 계층 conformal 구간\cite{Dunn2023}은 학습에서 뺀 관측의 평균 0.86을 포함하였다(95\% CI [0.80, 0.92], 명목 0.90, 4지역). 보정 지역이 4개이므로 유한 표본 보장은 없다(앞선 규약으로 내부 사전 등록; 보충 표~S12).
알래스카에서 conformalized quantile regression으로 보정한 분위 구간\cite{romano2019_conformalized_quantile_regression}은 학습에서 뺀 네 지역에서 0.32--0.73을 포함하였다(보충 표~S12).
계층 구간의 폭을 정하는 다섯 방식에서 포함률은 0.86--0.95, 평균 폭은 81.3--132.4~cm였다(그림~6d).
구간 점수\cite{GneitingRaftery2007}에서는 셀별 흐름 척도와 지역 안에서 폭을 섞은 위약 사이, 그리고 어느 척도와 상수 폭 사이에서도 차이를 확인하지 못했다.
라벨 10개에서 계층 예측 분포와 보정된 상수 폭 구간 사이의 구간 점수 차이는 확인하지 못했다(+9.5~cm, Holm 보정 \textit{P} = 0.277).
라벨 40개와 160개에서 잔차 ML 구간은 레나델타, 캐나다, 알래스카에서 0.84--0.90을 포함하였다(그림~6e).
